% Simple example for RandAssign
% Geoffrey M. Poore, 2015
% Adapté en français par Mikael LE MENTEC 2021
% License:  Creative Commons Zero (CC0) license
%           https://creativecommons.org/about/cc0
%
% Manual compile: run LaTeX, then PythonTeX, then LaTeX again
% Generate randomized assignments:  run `randassign example_simple.tex`

\documentclass{article}

\usepackage{nopageno}
\usepackage{pythontex}
\usepackage{xcolor}
\usepackage{tikz}

\begin{pycode}
from randassign import RandAssign
ra = RandAssign()
from math import *
import random
\end{pycode}

\title{Exemple simple pour RandAssign \\ {\Large Attempt: \input{attempt.tex}}}
\author{NOM : \underline{~~~~\input{name.tex}~~~~}}

\begin{document}

\maketitle

\begin{pycode}
a = random.randint(1, 10)
b = random.randint(1, 10)
c = sqrt(a**2 + b**2)
# Record the answer as formatted LaTeX with its unit, the way every other
# assessment here does. A bare float reaches the key with no unit and no control
# over how it is rendered. The units are written as literal "~cm" to match the
# problem statement above; this document does not load siunitx, so \SI would be
# an undefined control sequence in the generated solutions.
ra.addsoln('{0:.2f}~cm'.format(c))
\end{pycode}

Un triangle rectangle a une base de \py{a}~cm et une hauteur de \py{b}~cm. Déterminez la longueur de l'hypoténuse.



\begin{pysub}
\begin{tikzpicture}
\draw[style=help lines] (0,0) grid (10,10);
\filldraw[fill=blue!40!white] (0,0) -- (!{a},0) -- (!{a},!{b}) -- cycle;
\end{tikzpicture}
\end{pysub}

\end{document}
